\chapter{Resonances and Dressed Propagators}
\label{ch:volume-ii-resonances-dressed-propagators}
\ConventionsLorentzian


\section{Above-Threshold Exchange}

Continue the two-scalar model of the preceding chapter,
\[
  \mathcal L
  =
  -\frac12\partial_\mu\phi_1\partial^\mu\phi_1
  -\frac12m_1^2\phi_1^2
  -\frac12\partial_\mu\phi_2\partial^\mu\phi_2
  -\frac12m_2^2\phi_2^2
  -\frac12g\phi_1^2\phi_2 ,
\]
with \(m_1,m_2>0\) and \(g\in\R\). Let
\[
  K=k_1+k_2,\qquad s=-K^2,
\]
where \(k_1,k_2\) are the incoming \(\phi_1\) momenta. The \(s\)-channel
tree contribution to the connected invariant amplitude is
\[
  \mathcal M_s^{\mathrm{tree}}(s)
  =
  \frac{g^2}{m_2^2-s-\ii0}.
\]

Assume now
\[
  m_2>2m_1 .
\]
The value \(s=m_2^2\) lies in the two-\(\phi_1\) scattering region
\(s\ge4m_1^2\). The relevant scattering problem is therefore a problem with
stable external \(\phi_1\) states and an open two-particle continuum in the
same channel as the exchanged \(\phi_2\) line. The exact physical-axis
\(S\)-matrix is a bounded unitary boundary value; the pole suggested by the
tree graph is the lowest-order expansion of an analytic structure reorganized
by the exact two-point function of \(\phi_2\).

\begin{center}
\begin{tikzpicture}[>=Stealth,scale=0.94,
  one/.style={line width=0.85pt},
  two/.style={dash pattern=on 3pt off 2pt,line width=0.95pt},
  cut/.style={red!75!black,line width=1.6pt}]
  \begin{scope}
    \draw[one,->] (-1.65,0.9)--(-0.45,0.18);
    \draw[one,->] (-1.65,-0.9)--(-0.45,-0.18);
    \draw[two] (-0.45,0)--(1.15,0);
    \draw[one,->] (1.15,0.18)--(2.35,0.9);
    \draw[one,->] (1.15,-0.18)--(2.35,-0.9);
    \node at (0.35,0.35) {\(\phi_2\)};
    \node at (0,-1.35) {\(K^2=-s\)};
  \end{scope}
  \begin{scope}[xshift=5.1cm]
    \draw[->,line width=0.75pt] (-0.35,0)--(4.4,0) node[right] {\(s\)};
    \draw[->,line width=0.75pt] (0,-1.15)--(0,1.75)
      node[above] {\(\operatorname{Im}s\)};
    \draw[cut] (1.55,0)--(4.1,0);
    \node[red!75!black,below] at (1.55,-0.12) {\(4m_1^2\)};
    \fill[orange!90!black] (2.85,0) circle (2.4pt);
    \node[orange!90!black,below] at (2.85,-0.18) {\(m_2^2\)};
    \node[red!75!black,align=center] at (3.72,1.05)
      {open\\continuum};
    \draw[red!75!black,->,line width=0.8pt] (3.72,0.68)--(3.62,0.11);
    \draw[orange!90!black,->,line width=0.8pt] (2.85,0.65)--(2.85,0.12);
  \end{scope}
\end{tikzpicture}
\end{center}

The partial-wave normalization used below is the ordered \(16\pi\) convention
for the labeled \(\phi_1\phi_1\to\phi_1\phi_1\) amplitude:
\begin{equation}
  \mathcal M(s,z)
  =
  16\pi\sum_{\ell=0}^\infty(2\ell+1)\mathcal M_\ell(s)P_\ell(z),
  \qquad z=\cos\theta,
  \label{eq:resonance-partial-wave-expansion}
\end{equation}
Passing to the unordered identical-boson scattering channel amounts to the Bose
symmetrization displayed in
Equation~\eqref{eq:identical-partial-wave-conversion}; this chapter tracks the
analytic motion of fixed ordered-channel partial waves, so
\eqref{eq:resonance-partial-wave-expansion} is kept throughout.
In an elastic one-channel interval,
\begin{equation}
  S_\ell(s)=1+2\ii\,\rho_1(s)\mathcal M_\ell(s),
  \qquad
  \rho_1(s)=\sqrt{1-\frac{4m_1^2}{s}} .
  \label{eq:resonance-elastic-s-matrix-normalization}
\end{equation}
The branch of \(\rho_1\) is chosen by \(\rho_1(s)>0\) for \(s>4m_1^2\)
on the physical upper boundary of the first sheet, in the terminology of
Section~\ref{sec:physical-region-and-first-sheet}. Elastic unitarity gives
\begin{equation}
  |S_\ell(s)|=1 .
  \label{eq:resonance-elastic-unitarity}
\end{equation}
With inelastic channels included on the same Hilbert space, the elastic
matrix element instead obeys \(|S_\ell(s)|\le1\). Thus the exact physical-axis
partial wave is finite. The divergence of
\(\mathcal M_s^{\mathrm{tree}}\) at \(s=m_2^2\) records that the perturbative
description has expanded around a stable bare line in a region where the
state couples to an open decay channel. The exact object is a resonance: a
nearby pole of the continuation of \(S_\ell(s)\), or equivalently of
\(\mathcal M_\ell(s)\), through the threshold cut. On the physical sheet the
same channel is represented by the continuum part of the spectral measure;
an isolated real mass atom occurs only for a stable one-particle state.

\section{Outgoing Poles in a One-Dimensional Model}
\label{sec:resonance-one-dimensional-outgoing-poles}

This section is the detailed one-dimensional model promised in
Section~\ref{sec:bound-state-one-dimensional-orientation}.  It extends the
bound-state pole criterion to outgoing resonance poles; the underlying
Hilbert-space interpretation of generalized scattering states is the one fixed
in Section~\ref{sec:rigged-hilbert-continuous-spectrum}.  Let
\[
  H=-\frac{1}{2\mu}\frac{\dd^2}{\dd x^2}+V(x),
\]
where \(V\) is real and short range. For \(x\to+\infty\), use the outgoing
solution
\[
  \psi_k^{\mathrm{out}}(x,t)
  \sim
  S(k)^{-1}\ee^{-\ii kx-\ii Et}
  +
  \ee^{\ii kx-\ii Et},
  \qquad
  E=\frac{k^2}{2\mu}.
\]
If \(S(k)\) has a pole at \(k=k_\ast\), then the incoming coefficient
vanishes at that analytically continued value and the asymptotic solution is
purely outgoing. Write
\[
  E_\ast=\frac{k_\ast^2}{2\mu}=\omega-\ii\gamma,
  \qquad
  \gamma>0.
\]
Then
\[
  \psi_{k_\ast}^{\mathrm{out}}(x,t)
  \sim
  \ee^{\ii k_\ast x-\ii\omega t-\gamma t}.
\]
The negative imaginary part of \(E_\ast\) is the exponential decay rate of
the time dependence. The same wave grows in space along the outgoing
direction. This spatial growth is forced by probability conservation for a
state with only outgoing flux: a uniformly decaying packet cannot also be a
normalizable stationary eigenvector of a self-adjoint Hamiltonian. The pole
therefore defines a singular outgoing functional, often called a resonance or
Gamow vector, associated with the analytically continued scattering function.

\begin{center}
\begin{tikzpicture}[>=Stealth,scale=0.9]
  \begin{scope}
    \draw[->,line width=0.7pt] (-0.15,0)--(5.65,0) node[right] {\(x\)};
    \draw[->,line width=0.7pt] (0,-1.65)--(0,2.15) node[above] {\(V(x)\)};
    \draw[blue!70!black,line width=1.15pt]
      (0.25,1.75) .. controls (0.55,-0.95) and (1.08,-1.5) .. (1.48,-0.18)
      .. controls (1.85,0.92) and (2.18,0.78) .. (2.55,0.12)
      .. controls (3.1,-0.65) and (3.85,0.2) .. (5.2,0.05);
    \draw[orange!90!black,line width=1.05pt]
      (1.05,0.36) .. controls (1.3,0.75) and (1.58,0.75) .. (1.82,0.28)
      .. controls (2.02,-0.12) and (2.32,0.04) .. (2.66,0.36)
      .. controls (3.08,0.7) and (3.72,0.28) .. (4.92,-0.18);
    \foreach \x in {1.35,1.82,2.42,3.42,4.25}
      \draw[brown!70!black,->,line width=0.8pt] (\x,0.18)--++(0,-0.48);
    \draw[red!80!black,->,line width=1.05pt] (4.05,0.92)--(5.05,0.92);
    \node[red!80!black,align=left,anchor=west,font=\small] at (5.16,0.92)
      {outgoing\\flux};
    \node[orange!90!black,align=center,font=\small] at (1.08,1.05)
      {temporarily\\localized};
  \end{scope}
  \begin{scope}[xshift=8.0cm]
    \draw[->,line width=0.7pt] (-1.65,0)--(1.95,0)
      node[right] {\(\operatorname{Re}k\)};
    \draw[->,line width=0.7pt] (0,-1.55)--(0,1.9)
      node[above] {\(\operatorname{Im}k\)};
    \draw[blue!70!black,line width=1.1pt,->] (0,0)--(1.45,0);
    \fill[red!75!black] (0,0.85) circle (2.3pt);
    \node[red!75!black,anchor=west,font=\small] at (0.18,0.86)
      {bound};
    \fill[orange!90!black] (0.92,-0.62) circle (2.4pt);
    \node[orange!90!black,anchor=west,font=\small] at (1.08,-0.62)
      {\(k_\ast\)};
    \draw[orange!90!black,->,line width=0.8pt] (0.55,-0.18)--(0.86,-0.55);
    \node at (0,-1.95) {\(E=k^2/(2\mu)\)};
  \end{scope}
\end{tikzpicture}
\end{center}

Because \(E=k^2/(2\mu)\), the energy plane is a two-sheeted cover of the
\(k\)-plane. Bound-state poles lie on the positive imaginary \(k\)-axis and
map to negative real energies on the physical sheet. Resonance poles with
\(\operatorname{Re}k_\ast>0\) and \(\operatorname{Im}k_\ast<0\) map to
complex energies below the positive real axis on the sheet reached by
continuing through the scattering cut.

\begin{center}
\begin{tikzpicture}[>=Stealth,scale=0.92]
  \draw[->,line width=0.7pt] (-2.25,0)--(4.45,0)
    node[right] {\(\operatorname{Re}E\)};
  \draw[->,line width=0.7pt] (0,-1.65)--(0,2.05)
    node[above] {\(\operatorname{Im}E\)};
  \draw[blue!70!black,line width=1.05pt,->] (0,0)--(3.85,0);
  \node[blue!70!black,below,font=\small] at (2.55,-0.12)
    {physical continuum};
  \fill[red!75!black] (-1.18,0) circle (2.4pt);
  \fill[red!75!black] (-0.62,0) circle (2.4pt);
  \node[red!75!black,align=center,font=\small] at (-1.22,0.72)
    {bound\\states};
  \fill[orange!90!black] (2.35,-0.78) circle (2.5pt);
  \node[orange!90!black,anchor=west,font=\small] at (2.52,-0.78)
    {resonance pole};
  \draw[orange!90!black,dashed,line width=0.75pt] (2.35,-0.78)--(2.35,0);
  \draw[purple,->,line width=0.95pt]
    (1.0,0.72) .. controls (1.75,1.12) and (2.62,0.72) .. (2.82,0.08);
  \draw[green!55!black,->,line width=0.95pt]
    (0.52,-0.12) .. controls (0.82,-0.82) and (1.62,-1.18) .. (2.22,-0.84);
  \node[purple,font=\small] at (2.26,1.03) {first sheet};
  \node[green!55!black,font=\small] at (1.24,-1.3) {second sheet};
  \fill (0,0) circle (1.8pt) node[below left] {\(0\)};
\end{tikzpicture}
\end{center}

\section{The Dressed Propagator}

Return to the relativistic exchange model. Define the exact time-ordered
two-point function of \(\phi_2\) in momentum space by
\[
  D_2(k)
  =
  \int \dd^4x\,\ee^{-\ii k\cdot x}
  \bra{\vac}T\phi_2(x)\phi_2(0)\ket{\vac} .
\]
Let \(\Sigma(k)\) be the sum of amputated one-particle-irreducible two-point
insertions for the \(\phi_2\) line, with the convention
\[
  D_2(k)=\frac{-\ii}{k^2+m_{2,\mathrm{bare}}^2-\ii0-\Sigma(k)} .
\]
This is the geometric series obtained by inserting 1PI two-point kernels into
the bare propagator.

\begin{center}
\begin{tikzpicture}[>=Stealth,scale=0.84,
  two/.style={dash pattern=on 3pt off 2pt,line width=0.85pt},
  one/.style={line width=0.85pt},
  blob/.style={draw,circle,minimum size=0.58cm,inner sep=0pt}]
  \node at (-5.25,0) {\(D_2=\)};
  \draw[two] (-4.45,0)--(-3.35,0);
  \node at (-2.9,0) {\(+\)};
  \draw[two] (-2.35,0)--(-1.62,0);
  \node[blob] at (-1.25,0) {\scriptsize 1PI};
  \draw[two] (-0.88,0)--(-0.15,0);
  \node at (0.35,0) {\(+\)};
  \draw[two] (0.9,0)--(1.45,0);
  \node[blob] at (1.82,0) {\scriptsize 1PI};
  \draw[two] (2.18,0)--(2.72,0);
  \node[blob] at (3.1,0) {\scriptsize 1PI};
  \draw[two] (3.47,0)--(4.02,0);
  \node at (4.55,0) {\(+\cdots\)};
  \begin{scope}[yshift=-2.1cm]
    \node at (-3.75,0) {\(\ii\Sigma_1(k)=\)};
    \draw[two] (-2.45,0)--(-1.75,0);
    \draw[two] (-0.55,0)--(0.15,0);
    \draw[one] (-1.15,0) circle (0.6);
    \node[blue!70!black] at (-1.12,0.88) {\(\ell\)};
    \node[blue!70!black] at (-1.08,-0.88) {\(k-\ell\)};
    \node at (0.82,0) {\(+O(g^4)\)};
  \end{scope}
\end{tikzpicture}
\end{center}

At one loop,
\[
\begin{aligned}
  \ii\Sigma_1(k)
  &=
  \frac{(-\ii g)^2}{2}
  \int\frac{\dd^D\ell}{(2\pi)^D}\,
  \frac{-\ii}{\ell^2+m_1^2-\ii0}\,
  \frac{-\ii}{(k-\ell)^2+m_1^2-\ii0}.
\end{aligned}
\]
The factor \(1/2\) is the symmetry factor for the two identical \(\phi_1\)
lines in the loop. Combining denominators gives
\[
  \ii\Sigma_1(k)
  =
  \frac{g^2}{2}
  \int_0^1\dd x
  \int\frac{\dd^D\ell}{(2\pi)^D}\,
  \frac{1}{
    \left[\ell^2(1-x)+(k-\ell)^2x+m_1^2-\ii0\right]^2}.
\]
Shift \(\ell^\mu\mapsto\ell^\mu+xk^\mu\), then rotate the shifted energy
contour according to the Feynman prescription. The Euclidean expression is
\[
  \ii\Sigma_1(k)
  =
  \ii\frac{g^2}{2}
  \int_0^1\dd x
  \int\frac{\dd^D\ell_E}{(2\pi)^D}\,
  \frac{1}{
    \left[\ell_E^2+k^2x(1-x)+m_1^2\right]^2}.
\]
The rotation is a contour statement, not a change of notation: for the
shifted loop variable the pole displacements prescribed by \(-\ii0\) leave a
sector through which the \(\ell^0\)-contour can be moved to the Euclidean
axis without crossing singularities.  After the contour has been rotated, the
Euclidean denominator is an ordinary real denominator.  Its boundary value on a
timelike cut is obtained by analytically continuing the resulting function with
the original Feynman prescription; it is not encoded by leaving a residual
\(-\ii0\) in the Euclidean integrand.

\begin{center}
\begin{tikzpicture}[>=Stealth,scale=0.9]
  \draw[->,line width=0.7pt] (-2.8,0)--(2.8,0)
    node[right] {\(\operatorname{Re}\ell^0\)};
  \draw[->,line width=0.7pt] (0,-2.0)--(0,2.2)
    node[above] {\(\operatorname{Im}\ell^0\)};
  \draw[purple,line width=1.0pt,->] (-2.45,0)--(2.25,0);
  \draw[green!55!black,line width=1.0pt,->] (0,-1.7)--(0,1.75);
  \fill[red!75!black] (-1.12,0.35) circle (2.3pt);
  \fill[red!75!black] (1.15,-0.35) circle (2.3pt);
  \node[red!75!black,font=\small] at (-1.62,0.82)
    {Feynman poles};
  \draw[orange!90!black,->,line width=1.0pt]
    (1.6,0.18) .. controls (1.25,0.92) and (0.62,1.28) .. (0.08,1.34);
  \node[purple,font=\small,below] at (-1.62,-0.12) {Lorentzian};
  \node[green!55!black,font=\small,anchor=west] at (0.18,1.58)
    {Euclidean};
\end{tikzpicture}
\end{center}

In four dimensions, with a Euclidean cutoff \(|\ell_E|<\Lambda\) and with
\(k\) in the domain reached by the contour rotation,
\[
  \Sigma_1(k)
  =
  \frac{g^2}{32\pi^2}
  \int_0^1\dd x
  \left[
    \log
    \frac{\Lambda^2}{k^2x(1-x)+m_1^2}
    -1+O(\Lambda^{-2})
  \right].
\]
The divergent part is independent of \(k\). Introduce
\[
  m_{2,\mathrm{bare}}^2=(m_2^R)^2+\delta m_2^2,
  \qquad
  \Sigma_1(k)=\Sigma_1^R(k)+\delta m_2^2 ,
\]
where the counterterm cancels the cutoff-dependent constant. One convenient
subtraction is
\[
  \Sigma_1^R(k)
  =
  -\frac{g^2}{32\pi^2}
  \int_0^1\dd x\,
  \log
  \frac{k^2x(1-x)+m_1^2}{m_1^2}.
\]
This is first an analytic function on the Euclidean side of the cut.  The
physical boundary value is specified only after the external invariant is
continued to a timelike value.

\section{The Threshold Cut and the Cut Width}

For the \(s\)-channel exchange, \(k=K=k_1+k_2\) and \(K^2=-s\). The dressed
contribution is
\[
  \mathcal M_s(s)
  =
  \frac{g^2}{F(s)},
  \qquad
  F(s)=
  -s+(m_2^R)^2-\ii0-\Sigma_1^R(K).
\]
Equivalently,
\[
  F(s)
  =
  -s+(m_2^R)^2-\ii0
  +
  \frac{g^2}{32\pi^2}
  \int_0^1\dd x\,
  \log\left[
    1-\frac{s\,x(1-x)}{m_1^2}-\ii0
  \right].
\]
Since
\[
  0\le x(1-x)\le\frac14,
\]
the logarithm develops a discontinuity precisely for
\[
  s\ge4m_1^2 .
\]
For real \(s>4m_1^2\), define
\[
  \rho_1(s)=\sqrt{1-\frac{4m_1^2}{s}},
  \qquad
  x_\pm=\frac12\left(1\pm\rho_1(s)\right).
\]
The argument of the logarithm is negative for \(x\in[x_-,x_+]\). On the
physical upper boundary the prescription gives
\[
  \log(-|a|-\ii0)=\log|a|-\ii\pi ,
\]
and therefore
\[
  \operatorname{Im}\Sigma_1^R(K)\big|_{s+\ii0}
  =
  \frac{g^2}{32\pi}\rho_1(s).
\]
This boundary value is fixed by the same Wick-rotation prescription that
relates the Euclidean Green function to the Lorentzian time-ordered Green
function. For real positive energy above the two-particle threshold,
\(k^0\) is reached by analytic continuation from the imaginary \(k^0\)-axis
to the positive real axis through the upper half-plane.

\begin{center}
\begin{tikzpicture}[>=Stealth,scale=0.92]
  \begin{scope}
    \draw[->,line width=0.7pt] (-2.55,0)--(3.0,0)
      node[right] {\(\operatorname{Re}k^0\)};
    \draw[->,line width=0.7pt] (0,-2.0)--(0,2.25)
      node[above] {\(\operatorname{Im}k^0\)};
    \draw[green!55!black,line width=1.0pt] (0,-1.72)--(0,1.85);
    \draw[red!75!black,decorate,
      decoration={snake,amplitude=0.55mm,segment length=2.7mm},
      line width=0.95pt] (1.05,0)--(2.75,0);
    \draw[red!75!black,decorate,
      decoration={snake,amplitude=0.55mm,segment length=2.7mm},
      line width=0.95pt] (-2.35,0)--(-1.05,0);
    \draw[orange!90!black,->,line width=1.05pt]
      (0.22,1.35) .. controls (1.18,1.2) and (1.88,0.72) .. (2.12,0.12);
    \node[orange!90!black,align=center,font=\small] at (1.65,1.48)
      {continue\\from above};
  \end{scope}
  \begin{scope}[xshift=6.6cm]
    \draw[->,line width=0.7pt] (-0.4,0)--(5.05,0) node[right] {\(s\)};
    \draw[->,line width=0.7pt] (0,-1.75)--(0,2.0)
      node[above] {\(\operatorname{Im}s\)};
    \fill (1.28,0) circle (2pt) node[below] {\(0\)};
    \fill (2.62,0) circle (2pt) node[below] {\(4m_1^2\)};
    \draw[red!75!black,decorate,
      decoration={snake,amplitude=0.55mm,segment length=2.7mm},
      line width=1.05pt] (2.62,0)--(4.78,0);
    \node[purple,align=center,font=\small] at (3.72,1.1)
      {\(\operatorname{Im}\log=-\pi\)\\upper edge};
    \node[cyan!70!black,align=center,font=\small] at (3.72,-1.08)
      {\(\operatorname{Im}\log=+\pi\)\\lower edge};
    \draw[purple,->,line width=0.9pt] (3.2,0.78)--(3.2,0.12);
    \draw[cyan!70!black,->,line width=0.9pt] (3.2,-0.78)--(3.2,-0.12);
  \end{scope}
\end{tikzpicture}
\end{center}

Choose a real resonance mass parameter \(M_R\) by the condition that the real
part of the denominator vanishes at \(s=M_R^2\),
\[
  \operatorname{Re}F(M_R^2+\ii0)=0 .
\]
At weak coupling, and to leading order in the variation of the real part of
\(\Sigma_1^R\), the amplitude near this point is
\[
  \mathcal M_s(s)
  \simeq
  \frac{g^2}{M_R^2-s-\ii M_R\Gamma_{\rm cut}(M_R)},
\]
where
\[
  M_R\Gamma_{\rm cut}(M_R)
  =
  \frac{g^2}{32\pi}
  \sqrt{1-\frac{4m_1^2}{M_R^2}}
  +O(g^4).
\]
The subscript ``cut'' is deliberate.  The number above is extracted from the
absorptive part of the two-point function on the physical cut.  It is the
quantity computed by the inclusive cut, or Fermi golden rule, calculation at
the real energy \(M_R\).  It is not, by itself, a definition of every
resonance width.

\section{Width Parameters and Their Comparison}

There are several operationally different quantities that are often denoted by
the same letter \(\Gamma\).  In a narrow isolated resonance they agree to the
order at which the local Breit--Wigner approximation is meaningful, but their
definitions are distinct.

First, the cut width \(\Gamma_{\rm cut}(M)\) is defined at a real energy
\(M\) by the inclusive transition rate obtained from the discontinuity of the
self-energy.  For a scalar resonance coupled to stable final channels \(X\),
the perturbative expression has the form
\[
  \Gamma_{\rm cut}(M)
  =
  \frac{1}{2M}
  \sum_X
  \int \dd\Phi_X(P)\,
  \left|\mathcal A(R(P)\to X)\right|^2 ,
  \qquad P^2=-M^2,\quad P^0>0,
\]
with the symmetry factors included in the invariant phase-space measure
\(\dd\Phi_X(P)\) and in the amplitude \(\mathcal A\).  This formula should be
read as a weak-coupling construction using the stable state \(R\) of the
decoupled or lower-order theory as an interpolating state.  Once the channel is
open, the exact theory has stable asymptotic states in the final channels, not
an exact external one-resonance state.

Second, the line-shape width \(\Gamma_{\rm line}\) is an observable extracted
from a physical inclusive cross section or event distribution.  After a
specified smooth background subtraction and a specified resolution function,
one may define it as the full width at half maximum in the energy variable
\(E=\sqrt{s}\) of the resonant contribution.  This definition depends on the
chosen process, on the observable used to form the inclusive distribution, and
on how nonresonant background is separated from the pole term.  It is therefore
not an invariant analytic datum for a broad resonance.

Third, the Gamow time width, or complex-energy pole width, is defined from the
second-sheet pole after choosing an energy square root.  If \(s_\ast\) is a
simple resonance pole and the square-root branch is chosen so that
\(\operatorname{Re}\sqrt{s_\ast}>0\) near the resonance, write
\[
  E_\ast=\sqrt{s_\ast}
  =
  M_{\rm pole}-\frac{\ii}{2}\Gamma_{\rm pole},
  \qquad M_{\rm pole}>0,\quad \Gamma_{\rm pole}>0.
\]
Equivalently,
\[
  s_\ast
  =
  M_{\rm pole}^2
  -\ii M_{\rm pole}\Gamma_{\rm pole}
  -\frac14\Gamma_{\rm pole}^2 .
\]
This definition is independent of a particular production process once the
pole and sheet have been specified.  The associated outgoing Gamow functional
evolves as
\[
  \ee^{-\ii E_\ast t}
  =
  \ee^{-\ii M_{\rm pole}t}\,
  \ee^{-\Gamma_{\rm pole}t/2},
\]
so its probability norm, in the formal Gamow description, decays as
\(\ee^{-\Gamma_{\rm pole}t}\).  This is the time-decay width of the pole
functional.  The statement concerns the pole contribution in the rigged-Hilbert
analytic continuation; an exact normalizable state in a self-adjoint
Hamiltonian does not have a purely exponential survival law for all times.

Fourth, one can define an \(s\)-plane pole-imaginary width by
\[
  M_s=\sqrt{\operatorname{Re}s_\ast},
  \qquad
  \Gamma_s=-\frac{\operatorname{Im}s_\ast}{M_s}.
\]
This convention uses the imaginary part of the pole in the \(s\)-plane rather
than the imaginary part of the energy \(E=\sqrt{s}\).  It agrees with
\(\Gamma_{\rm pole}\) only up to the kinematic conversion
\[
  \Gamma_s
  =
  \Gamma_{\rm pole}
  \left[
    1+O\!\left(\frac{\Gamma_{\rm pole}^2}{M_{\rm pole}^2}\right)
  \right].
\]

\begin{controlledapproximation}[Narrow-resonance comparison of widths]
Suppose that, in the energy window
\(|E-M_{\rm pole}|=O(\Gamma_{\rm pole})\), the continued stable-channel
amplitude is dominated by a single simple pole, the residue and the
nonresonant background are holomorphic and vary only on a scale of order
\(M_{\rm pole}\), and the pole is separated from thresholds and other
singularities by a squared-mass distance much larger than
\(M_{\rm pole}\Gamma_{\rm pole}\).  Then the ideal resonant line shape
obtained from the pole term, with the smooth residue replaced by its value at
the pole and with the smooth background removed, has
\[
  \Gamma_{\rm line}
  =
  \Gamma_{\rm pole}
  \left[
    1+O\!\left(\frac{\Gamma_{\rm pole}^2}{M_{\rm pole}^2}\right)
  \right],
\]
as follows.  Keeping only the pole term and writing
\(E=M_{\rm pole}+x\),
\[
  |s_\ast-E^2|^2
  =
  \left(
    2M_{\rm pole}x+x^2+\frac14\Gamma_{\rm pole}^2
  \right)^2
  +
  M_{\rm pole}^2\Gamma_{\rm pole}^2 .
\]
The half-maximum points of the pole term, measured in the energy variable, are
therefore
\[
  x_\pm
  =
  \pm\frac12\Gamma_{\rm pole}
  -\frac{1}{4M_{\rm pole}}\Gamma_{\rm pole}^2
  +O\!\left(\frac{\Gamma_{\rm pole}^3}{M_{\rm pole}^2}\right),
\]
so \(x_+-x_-=\Gamma_{\rm pole}[1+O(\Gamma_{\rm pole}^2/M_{\rm pole}^2)]\).
The common shift of the two half-maximum points records that the peak in
\(E\) is not centered exactly at \(M_{\rm pole}\) once the quadratic relation
\(s=E^2\) is retained.  If the residue, phase space, detector resolution, or
subtracted background varies appreciably over the same energy interval, the
measured line-shape width receives additional process-dependent corrections.
That is precisely why \(\Gamma_{\rm line}\) is not an invariant analytic datum
for a broad resonance.

The cut calculation at the real resonance energy gives
\[
  M_{\rm pole}\Gamma_{\rm pole}
  =
  M_R\Gamma_{\rm cut}(M_R)
  +O(g^4)
\]
in the weak-coupling scalar model above.  Indeed, on the second sheet the
discontinuity computation gives
\[
  F_{\rm II}(M_R^2)
  =
  -\ii M_R\Gamma_{\rm cut}(M_R)+O(g^4),
  \qquad
  F_{\rm II}'(M_R^2)=-1+O(g^2).
\]
Solving \(F_{\rm II}(s_\ast)=0\) by one Newton step therefore gives
\[
  s_\ast-M_R^2
  =
  -\ii M_R\Gamma_{\rm cut}(M_R)+O(g^4).
\]
Thus the common symbol \(\Gamma_R\) is justified only after these hypotheses
and approximations have been stated.  For a broad resonance, for an
overlapping pair of poles, or near a threshold branch point,
\(\Gamma_{\rm line}\), \(\Gamma_{\rm cut}\), \(\Gamma_{\rm pole}\), and
\(\Gamma_s\) need not coincide.
\end{controlledapproximation}

\section{Elastic Phase Motion}

In the same approximation, only the \(\ell=0\) partial wave carries the
resonant scalar exchange. Using
\[
  \mathcal M_s(s)
  \approx
  -16\pi\ii\sqrt{\frac{s}{s-4m_1^2}}\,
  \bigl(S_0(s)-1\bigr),
\]
one obtains, near the resonance,
\[
  S_0(s)
  \approx
  \frac{
    M_R^2-s+\ii W(s)
  }{
    M_R^2-s-\ii W(s)
  },
  \qquad
  W(s)
  =
  \frac{g^2}{32\pi}\rho_1(s).
\]
For real \(s\) in the elastic interval, the numerator is the complex
conjugate of the denominator, so \(|S_0(s)|=1\). The resonance is expressed
as rapid phase motion. If
\[
  S_0(s)=\ee^{2\ii\delta_0(s)},
\]
then the phase shift \(\delta_0\) changes rapidly through the point where
\(\operatorname{Re}F(s)\) vanishes.

\begin{center}
\begin{tikzpicture}[>=Stealth,scale=0.95]
  \begin{scope}
    \draw[->,line width=0.75pt] (-2.0,0)--(2.15,0)
      node[right] {\(\operatorname{Re}S_0\)};
    \draw[->,line width=0.75pt] (0,-2.0)--(0,2.15)
      node[above] {\(\operatorname{Im}S_0\)};
    \draw[green!55!black,dashed,line width=0.9pt] (0,0) circle (1.42);
    \draw[blue!70!black,line width=1.25pt]
      (1.38,-0.08) arc[start angle=-3,end angle=345,radius=1.42];
    \draw[blue!70!black,->,line width=1.0pt]
      (0.2,-1.4) arc[start angle=-82,end angle=-20,radius=1.42];
    \node[green!55!black,anchor=west,font=\small] at (1.78,1.0)
      {\(|S_0|=1\)};
  \end{scope}
  \begin{scope}[xshift=6.4cm]
    \draw[->,line width=0.75pt] (-0.25,0)--(5.2,0) node[right] {\(s\)};
    \draw[->,line width=0.75pt] (0,-1.45)--(0,1.9)
      node[above] {\(\operatorname{Im}s\)};
    \draw[red!75!black,decorate,
      decoration={snake,amplitude=0.55mm,segment length=2.7mm},
      line width=1.05pt] (2.15,0)--(4.9,0);
    \node[red!75!black,below] at (2.15,-0.12) {\(4m_1^2\)};
    \fill[orange!90!black] (3.35,-0.72) circle (2.4pt);
    \node[orange!90!black,anchor=west,font=\small] at (3.52,-0.72)
      {second-sheet pole};
    \draw[orange!90!black,dashed,line width=0.75pt] (3.35,-0.72)--(3.35,0);
    \node[orange!90!black,above,font=\small] at (3.35,0.2) {\(M_R^2\)};
  \end{scope}
\end{tikzpicture}
\end{center}

\section{The Second Sheet}

The physical amplitude is the boundary value reached from the upper half of
the first sheet. Across the threshold cut the square-root branch changes
sign. Equivalently, the logarithm in \(F(s)\) has the discontinuity
\[
  \log(-|a|-\ii0)-\log(-|a|+\ii0)=-2\pi\ii .
\]
The first sheet contains the physical branch cut beginning at \(4m_1^2\).
The resonance pole is reached by continuing through this cut to the adjacent
sheet.

The first-sheet denominator has a useful sign property. Continue \(F(s)\)
from the real interval \(s<4m_1^2\) into the upper half-plane. For
\(\operatorname{Im}s>0\),
\[
  \operatorname{Im}\!\left[
    \log\left(1-\frac{s\,x(1-x)}{m_1^2}\right)
  \right]<0,
  \qquad
  \operatorname{Im}(-s)<0,
\]
for \(0<x<1\). Hence \(\operatorname{Im}F(s)<0\) on the upper half of the
first sheet, so \(F(s)\) has no zero there. The lower half-plane gives the
conjugate sign statement. Thus the resonance associated with the dressed
\(\phi_2\) exchange is not a pole of either half of the first sheet; its pole
appears only after passing through the cut.

Near a narrow resonance, this continuation through the \(\phi_1\phi_1\)
unitarity cut places the pole on the second sheet relative to that cut at
\[
  s_\ast
  =
  M_R^2-\ii M_R\Gamma_{\rm cut}(M_R)+O(g^4),
\]
with \(\Gamma_{\rm cut}(M_R)>0\).  Equivalently, the pole width defined by
\(E_\ast=\sqrt{s_\ast}=M_{\rm pole}-\ii\Gamma_{\rm pole}/2\) obeys
\[
  M_{\rm pole}\Gamma_{\rm pole}
  =
  M_R\Gamma_{\rm cut}(M_R)+O(g^4).
\]
More invariantly, \(s_\ast\) is a zero of the second-sheet continuation of
\(F(s)\):
\[
  F_{\mathrm{II}}(s_\ast)=0 .
\]
The number \(s_\ast\) is an analytic datum of the stable-particle scattering
amplitude. Its position sets the resonance location and the pole lifetime
scale in the narrow-width regime; it is not the same kind of object as a
physical-axis line-shape width unless the narrow-resonance hypotheses of the
preceding section hold.

The same distinction is visible in the spectral representation. A stable
one-particle state gives an atom in the spectral measure and supplies a
Haag--Ruelle asymptotic sector. An open channel gives continuous spectral
measure beginning at threshold. A narrow resonance is a pronounced analytic
feature of that continuum: it can appear as a peak in the spectral density and
as a pole after analytic continuation, while the external scattering states
remain the stable states used in the construction of the \(S\)-matrix.

\section{Unstable Resonances and External-Leg Reduction}

The LSZ theorem applies to an external particle only when the theory supplies
an isolated real mass shell in the physical Hilbert space.  In the notation of
Chapter~\ref{ch:lsz-reduction}, this means an orthogonal projection
\[
  P_1=E(\Sigma_m^+)
\]
onto a Wigner one-particle subspace, a positive pole residue in a local-field
two-point function, and Haag--Ruelle incoming and outgoing states built from
that subspace.  A resonance pole
\(s_\ast=M_{\rm pole}^2-\ii M_{\rm pole}\Gamma_{\rm pole}
-\Gamma_{\rm pole}^2/4\) has
none of these properties.  It is a pole of an analytically continued stable
scattering amplitude on a nonphysical sheet; the associated outgoing Gamow
functional is not a normalizable vector in \(\Hilb\), and there is no
projection-valued spectral atom \(E(\{p:p^2=-s_\ast\})\).

Therefore a residue at \(s=s_\ast\) is not an LSZ external-leg residue.
Physical \(S\)-matrix elements with a resonance are matrix elements between
stable asymptotic states.  Near an isolated resonance pole the meromorphic
continuation of such a stable-particle amplitude may have the local form
\[
  \mathcal M_{a\to b}(s,\alpha,\beta)
  =
  \frac{r_{a}(\alpha)\,r_{b}(\beta)}{s_\ast-s}
  +
  \mathcal M_{\rm reg}(s,\alpha,\beta),
\]
where \(a,b\) label stable channels and \(\alpha,\beta\) denote the remaining
kinematic variables.  The functions \(r_a,r_b\) are pole residues of the
continued stable-channel amplitude.  They can be used to describe production
and decay through the resonance, but they do not define an incoming or
outgoing one-resonance Hilbert-space state.

Perturbation theory often packages this pole information in a
Stueckelberg--Veltman, or complex-mass-shell, prescription.  One analytically
continues the inverse dressed propagator to the pole,
\[
  F_{\rm II}(s_\ast)=0,
  \qquad
  Z_\ast^{-1}= -F_{\rm II}'(s_\ast),
\]
and replaces the resonant propagator near the pole by
\[
  \frac{g^2}{F_{\rm II}(s)}
  =
  \frac{g^2 Z_\ast}{s_\ast-s}
  +
  \text{terms holomorphic near }s_\ast .
\]
This is a pole-scheme parametrization of the stable-channel amplitude, not a
new LSZ theorem.  Its usefulness requires a controlled narrow-resonance
regime: the pole must be isolated from other nearby poles and branch points,
the width scale \(M_{\rm pole}\Gamma_{\rm pole}\) must be small compared with
the squared-mass distance to thresholds or singularities whose variation is
being neglected, and the nonresonant background must be slowly varying on the
same scale.  If the pole lies close to a threshold, or if several channels
produce nearby singularities, the simple Breit--Wigner or complex-mass-shell
description is replaced by the corresponding multi-channel analytic line
shape.  The external states remain the stable states in all cases.

\begin{openproblem}[External resonances]
There is a further problem, distinct from ordinary LSZ reduction: formulate a
theory of amplitudes whose external objects are unstable resonances.  Such a
theory cannot use normalizable one-resonance vectors in the physical Hilbert
space.  A satisfactory formulation would have to start from stable-channel
unitarity and locality, continue amplitudes to the relevant sheets, identify
field-redefinition-invariant pole data, and state factorization, covariance,
gauge-invariance, and threshold conditions for the resulting resonance
objects.  The complex-mass-shell prescription is a perturbative pole
parametrization; it does not by itself solve this nonperturbative
scattering-theoretic problem.
\end{openproblem}
